\section{Alzheimer's: desde la reproducción de mecanismos hasta un nodo prospectivo}

El caso de Alzheimer es más complejo que las reproducciones de resultados ocultos anteriores: el sistema primero reconstruye una cadena de mecanismos multietapa ACE--bradiquinina--B2R, y luego equipos colaboradores realizan experimentos prospectivos en uno de sus nodos. Por ello, incluye simultáneamente contexto público, recapitulación retrospectiva y validación prospectiva no publicada, y debe narrarse desglosando cada parte.

\subsection{Indicar primero el estado de colaboración y publicación}

El curso colabora con investigadores del MGH para reproducir la ruta de experimentación humana de años de duración y predecir independientemente en días el eje ACE--B2R patógeno. La diapositiva también indica datos no publicados, que necesitan revisión por pares. Aquí el «case study» no se refiere a un caso clínico ni a una recomendación terapéutica publicada, sino a la exhibición de una hipótesis de mecanismo y un flujo experimental.

\lecturefigure{slide-016.jpg}{Objetivos del caso Alzheimer y límites de datos no publicados}{Grabación oficial Stanford CS25 V6 Lecture 07 00:53:54}

\begin{warningbox}{Tres tipos de evidencia deben presentarse en columnas separadas}
La literatura pública ya ha discutido la relación entre bradiquinina, daño cognitivo y Alzheimer; el co-científico de IA reconstruye la cadena de mecanismos desde inhibidores de ACE hasta B2R; y equipos colaboradores realizaron experimentos prospectivos sobre el nodo B2R como se describió en clase. Los dos primeros no pueden probar al tercero, y antes de que sea publicado, el tercero tampoco puede escribirse como un mecanismo terapéutico confirmado.
\end{warningbox}

\subsection{Cómo se forma la cadena ACE--bradiquinina--B2R}

Los inhibidores de ACE se usan para tratar la hipertensión e inhiben la enzima convertidora de angiotensina. Además de participar en el sistema renina--angiotensina, ACE degrada bradiquinina; inhibir ACE podría aumentar la señalización de bradiquinina. El curso propone que la activación de B2R por bradiquinina desencadena una cascada inflamatoria que podría explicar la asociación observada entre los inhibidores de ACE y el riesgo de Alzheimer en ciertos grupos poblacionales.

\lecturefigure{slide-017.jpg}{Reproducción por IA de la cadena de mecanismos: inhibidor de ACE, bradiquinina y B2R}{Grabación oficial Stanford CS25 V6 Lecture 07 00:54:48}

\begin{knowledgebox}{Lectura de la figura: la cadena de cuatro pasos requiere verificación gradual}
A la izquierda se listan paradojas clínicas: los inhibidores de ACE son frecuentes en hipertensión, pero ciertos estudios observacionales los asocian con riesgo de Alzheimer; a la derecha está el mecanismo construido por IA. Se deben preguntar por separado: ¿el fármaco altera los niveles de bradiquinina? ¿La bradiquinina actúa sobre las células relevantes a través del receptor B2R? ¿El B2R causa fenotipos patológicos? ¿Bloquear el B2R puede revertir la condición? Si falla un solo eslabón, la narrativa causal completa debe revisarse.
\end{knowledgebox}

Se puede escribir la cadena de mecanismos como una hipótesis causal estructurada:
$$
\text{Inhibición de ACE}
\rightarrow \uparrow\text{bradiquinina}
\rightarrow \text{Activación de B2R}
\rightarrow \text{Fenotipo inflamatorio}.
$$
Donde:
\begin{itemize}
  \item La primera flecha depende del papel degradador de ACE sobre la bradiquinina;
  \item La segunda flecha representa una relación ligando--receptor;
  \item La tercera flecha conecta con neuroinflamación o fenotipos patológicos;
  \item La correlación observacional no determina automáticamente la dirección de las flechas ni descarta causas comunes y sesgos en el uso de medicamentos.
\end{itemize}

\begin{warningbox}{La fluidez de una historia causal puede generar sobreconfianza}
Los LLM son muy hábiles para tejer múltiples hechos locales reales en una cadena completa, pero la evidencia de cada arista puede provenir de especies, tejidos o condiciones experimentales diferentes. Un Agente científico debe producir procedencia a nivel de arista y experimentos refutables, en lugar de solo ofrecer una narrativa de mecanismo fluida.
\end{warningbox}

\subsection{Validación prospectiva: por qué un nodo es importante}

En clase se afirmó que el sistema identificó un nodo regulatorio ignorado y propuso usar antagonistas para intervenir sobre B2R. Equipos colaboradores probaron las predicciones en laboratorio e informaron que un nuevo inhibidor logró revertir el fenotipo asociado a la enfermedad. Si el candidato fue congelado antes del experimento, esto es más fuerte que una explicación retrospectiva porque el sistema asumió el riesgo de predicción fallible sobre los resultados.

\lecturefigure{slide-018.jpg}{Tres fases: predicción, experimentación y validación del nodo B2R}{Grabación oficial Stanford CS25 V6 Lecture 07 00:55:26; datos no publicados de clase}

\begin{knowledgebox}{Lectura de la figura: lo clave de prospectivo es el orden temporal}
La primera columna, Predicción, debe ocurrir antes de ver los resultados del experimento; la segunda columna, Experimentación, convierte el antagonista de B2R en una intervención ejecutable; la tercera columna, Validación, reporta la reversión del fenotipo. Si estos tres pasos tienen marcas de tiempo, informes congelados y métricas predefinidas, se reduce la selección a posteriori. La figura aún no muestra métodos completos ni repeticiones independientes, por lo que es más fuerte que una recapitulación pero más débil que evidencia pública revisable.
\end{knowledgebox}

Si con evidencia previa al experimento $D_0$ se forma la hipótesis $h$, y los resultados del experimento son $e$, la actualización puede escribirse como
$$
p(h\mid D_0,e)
=\frac{p(e\mid h,D_0)p(h\mid D_0)}
{p(e\mid D_0)}.
$$
Donde:
\begin{itemize}
  \item $p(h\mid D_0)$ es la credibilidad previa al experimento;
  \item $p(e\mid h,D_0)$ es la probabilidad de observar el resultado $e$ si el mecanismo se cumple;
  \item $p(e\mid D_0)$ incluye la posibilidad de que otros mecanismos produzcan el mismo resultado;
  \item Un éxito único aumenta la credibilidad, pero no permite empujar directamente el posterior a 1.
\end{itemize}

\teachervoice{00:53:33--00:56:04, el docente divide este caso en una recapitulación de un pipeline humano multiplataforma y una validación prospectiva del nodo B2R. Ambos aportes no deben mezclarse en una sola frase tipo «IA descubre la causa».}

\subsection{Cómo leer la comparación con base LLM}

La tabla del curso compara la cobertura de varios elementos de mecanismo entre el co-científico de IA y Claude 4 / GPT-5: acumulación inicial de bradiquinina, sensibilización específica a B2R asociada a APOE4, crisis energética neuronal, quinasas específicas y complejidad causal final. La tabla pretende mostrar que los sistemas multiagente pueden formar cadenas de mecanismos más completas, pero sigue siendo un análisis del contenido de salida, no un ensayo controlado aleatorizado sobre eficacia terapéutica.

\lecturefigure{slide-019.jpg}{Comparación de elementos de mecanismo entre co-científico de IA y base LLM}{Grabación oficial Stanford CS25 V6 Lecture 07 00:56:04}

\begin{knowledgebox}{Lectura de la figura: se compara «cobertura de mecanismos», no un ranking general de capacidad médica}
Cada fila corresponde a elementos de mecanismo que los investigadores consideran posteriormente, y las columnas registran si el sistema los acertó u omitió. Las ventajas del co-científico de IA podrían provenir de búsquedas más largas, agentes especializados, herramientas y torneos, no necesariamente porque el modelo base entienda mejor la medicina. Si las filas de evaluación las definen quienes ya vieron la respuesta objetivo, hay que vigilar fugas en el rúbrica. El uso más robusto es hacer ablaciones: quitar uno por uno los agentes, herramientas o cómputo para ver qué mecanismos desaparecen con ellos.
\end{knowledgebox}

\begin{warningbox}{Tres limitaciones de la tabla de benchmark}
Muestra una muestra pequeña; las filas de mecanismo pueden haber sido seleccionadas manualmente; «acertar una palabra» no es lo mismo que «proporcionar un modelo causal verificable». Esta tabla apoya hipótesis de diseño del sistema, pero no puede sostener por sí sola mecanismos de enfermedad ni decisiones clínicas.
\end{warningbox}

\subsection{Resumen de la sección}

El caso de Alzheimer ilustra un ciclo científico más completo: el contexto público ofrece pistas, los sistemas multiagente reconstruyen el mecanismo ACE--bradiquinina--B2R y equipos colaboradores realizan intervenciones prospectivas sobre B2R. Se acerca más al descubrimiento que las reproducciones de resultados ocultos, pero los datos de clase aún no están publicados. La forma correcta de leerlo es rastrear cada arista causal, congelar el tiempo de predicción y estratificar estrictamente benchmark, experimentación y evidencia clínica.
